\chapter{\abstractname}

With the increasing deployment of residential photovoltaic (PV) systems, battery energy storage systems, and electric vehicles (EVs), prosumers in low-voltage distribution networks are gaining both the need and the capability to manage their electricity consumption and generation flexibly. The energy scheduling decisions of multiple prosumers affect not only their individual operating costs but also the safety of the distribution network. Coordinated control under fluctuating electricity prices, EV mobility requirements, and grid safety constraints therefore represents an important challenge in building energy management.


This thesis investigates the coordinated scheduling of EV charging, stationary batteries, and PV generation for three prosumers connected to a low-voltage distribution network. A distribution network simulation environment is developed using pandapower and incorporates physical models of stationary batteries, EVs, PV systems, and distribution network safety constraints. Multi-agent deep reinforcement learning (MADRL) is employed to learn price-aware and resource-coupled control policies. Action feasibility checks, EV departure constraint handling, and an optional grid safety projection are applied before the actions are executed. The MADRL variants are compared with Local model predictive control (Local MPC), alternating direction method of multipliers-based MPC (ADMM-MPC), and a rule-based battery controller.


The controllers are evaluated over 15 days for three prosumers, resulting in 45 EV departure records. Among the investigated MADRL variants, Projected Hard provides the best overall balance between economic performance, EV feasibility, and distribution network safety. Its total cost is 186.04 EUR, which is 32.1\% lower than the 274.01 EUR achieved by the rule-based controller. However, it remains 188.90 EUR and 222.41 EUR more expensive than ADMM-MPC Hard and Local MPC Hard, respectively. Projected Hard reduces the number of voltage violations from three for the rule-based controller and 25 for Local MPC Hard to zero; ADMM-MPC Hard also records zero voltage violations, so that it can ensure grid safety. Moreover, Projected Hard satisfies the EV departure state-of-charge target in all 45 departure records, which behaves better compared with MPC methods.

The results show that MADRL can learn more flexible price-aware and resource-coupling policies than the rule-based controller. A progress penalty improves EV feasibility under soft constraints by replacing a sparse departure signal with denser temporal guidance, whereas hard constraint handling enforces EV charging feasibility at the action-execution layer. Grid safety projection eliminates voltage violations in the evaluated scenario, but modifies the actions selected by the learned policy and consequently reduces battery arbitrage performance. Overall, the results demonstrate the trade-off among economic cost, EV mobility requirements, and distribution network safety, and provide an experimental basis for coordinated EV--battery--PV control in low-voltage networks with multiple prosumers.
